% !TeX spellcheck = en_US
% !TeX program=pdflatex
% !BIB program=biber

%%%%%%%%%%%%%%%%%%%%%%%%%%%%%%%%%%%%%%%%%%%%%%%%%%%%%%%%%%%%%%%%%%%%%%%
% Options for the class:
% 'fullpaper' if you are submitting a 6-page paper or
% 'abstract' if you are submitting a 2-page extended abstract.
% 'review' (is the default) will anonymize the submission and will add the review line numers.
% 'final' will add a mandatory footer on the titlepage which must contain the correct volume. This will be provided once the paper is accepted and everything is ready for publication. This footer is only necessary for full papers.
%%%%%%%%%%%%%%%%%%%%%%%%%%%%%%%%%%%%%%%%%%%%%%%%%%%%%%%%%%%%%%%%%%%%%%%
\documentclass[final]{nldl}


%% Math
\usepackage{mathtools}

% Figures
\usepackage{graphicx}


%% Tables
\usepackage{booktabs}


%% Lists
\usepackage{enumitem}


% Algorithms
\usepackage{algorithm}
\usepackage{algorithmic}
\newcommand{\theHalgorithm}{\arabic{algorithm}}


%% Code
\usepackage{listings}
\lstset{
  basicstyle=\small\ttfamily,
  breaklines,
}

% References
\usepackage[backend=biber]{biblatex}
\addbibresource{references.bib}


% Links and hyperlinks
\usepackage{hyperref}
\usepackage{url}
\hypersetup{
  pdfusetitle,
  colorlinks,
  linkcolor = BrickRed,
  citecolor = NavyBlue,
  urlcolor  = Magenta!80!black,
}



% Replace with your title, authors and affiliation
\title{Recod.ai/LUC - Scientific Image Forgery Detection}
\author[1]{Max Bakke Seymour}
\author[1]{Einar Martin Søvik Bernsen}
\author[1]{Emre Balci}
\affil[1]{University of Bergen (UiB)}

\begin{document}
\maketitle

\begin{abstract}
TODO - "Short summary of the paper, including results. Often written last".
\end{abstract}

\section{Introduction}

\subsection{Copy-Move Forgery Detection}

Copy-Move Forgery is a class of image forgery where regions of an image are duplicated and moved to other areas of the image. Commonly, techniques such as rotation, flipping and scaling of the copied region are applied, as well as post-processing such as edge smoothing, brightness adjustment and adding extra noise. These techniques can lead to forgeries which are hard to detect both through manual review and automated detectors. This method of forgery can be used within scientific images in order to fabricate results, where the complexity of the figures creates additional difficulties. \cites{kaggle}
\subsection{Related works}
TODO: What have others tried to solve this problem? 
Are there similar problems studied?

\subsection{Objectives}
TODO
\\\\
We aimed to create a full pipeline for CMFD which could perform well on the Recoid.ai Scientific Image Forgery Detection dataset.


\subsection{Contributions}
TODO: What can the reader learn from this paper?

\section{Methods}

TODO: What did we do?
Reader should be able to reproduce results.

Basically, detailed breakdown of final implementation goes here.

\section{Results}

TODO: Show results of experiments. NB: Without interpretation at this stage!
Align with method section.

\section{Discussion}


TODO: Interpret results
Relate to objectives (did we achieve them?)
Compare against others (kaggle leaderboard?)
\printbibliography

\appendix
\section{Appendix Example}

\end{document}
